% P18-A1: Dung's arbitrary-framework semantic and graph-condition chain.
\chapter{Abstract Argumentation on Arbitrary Sets}
\label{sec:p18-arbitrary-argumentation}

The definitions of attack, defense and extension in
\cref{sec:p18-finite-argumentation} also apply when the argument set
is infinite. Finiteness enters the existence and computation
arguments, not those definitions. We now give the arbitrary-set
results and the graph conditions of Dung's original semantic
presentation \cite[sec.~2, pp.~325--334]{Dung1995}.

\begin{hbdefinition}{def:p18-arbitrary-framework}{The arbitrary-set framework and set-theoretic premises}
Let \(\mathcal F=(A,R)\), where \(A\) is any set, including the
empty set, and \(R\subseteq A\times A\). Retain exactly the
attack orientation, self-attacks, \(\operatorname{Att}(a)\),
\(S^+\), characteristic operator \(\mathsf F\), admissibility,
and the complete, preferred and stable definitions of
\cref{def:p18-finite-framework,def:p18-defense-operator,def:p18-extension-semantics}.
Write \(\operatorname{Adm}(\mathcal F)\) for the admissible sets.
The arrow \(a\to b\) abbreviates \((a,b)\in R\).

The arbitrary-set existence results below are proved in ordinary
set theory with the axiom of choice. We identify the uses:
Zorn's lemma supplies maximal admissible supersets; well-ordering
and transfinite recursion bound increasing constructions;
dependent choice converts a relation with an incoming edge
at every point of a nonempty set into an infinite backward
sequence. Finite specializations need none of these
infinite selection arguments.

A nonempty family of subsets is \emph{directed} when every
two members have a common containing member in that family.
A \emph{chain} is a family linearly ordered by inclusion.
A least upper bound and a greatest lower bound always
refer to the specified poset, not automatically to the
union and intersection in \(2^A\).
\end{hbdefinition}

\begin{hbtheorem}{thm:p18-arbitrary-admissibility}{proved=annex:p18.arbitrary.admissibility}{Defense, directed unions and preferred existence}
On every such framework, \(\mathsf F\) is monotone and maps
conflict-free sets to conflict-free sets. It maps admissible
sets to admissible sets. If \(S\) is admissible and
\(a,a'\in\mathsf F(S)\), then \(S\cup\{a\}\) is admissible
and still defends \(a'\).

The union of every nonempty directed family of admissible
sets is admissible and is its least upper bound in
\(\operatorname{Adm}(\mathcal F)\). The empty set is the
least admissible set. Every admissible set is contained
in a preferred extension, so preferred extensions exist.
In particular \(\operatorname{Adm}(\mathcal F)\) is a
pointed directed-complete partial order: it has a
least element and suprema of nonempty directed families.

Every stable extension is preferred, every preferred
extension is complete, and complete sets are exactly
the conflict-free fixed points of \(\mathsf F\).
Also \(S\) is stable exactly when \(S=A\setminus S^+\).
These assertions include empty frameworks and self-attacks
\cite[lemmas~10, 14--15, 18--19, 24; theorem~11]{Dung1995}.
\end{hbtheorem}

\section{Fixed points, limits and order structure}

\begin{hbdefinition}{def:p18-transfinite-closure}{Increasing iteration from an admissible seed}
For an admissible set \(S\), define the ordinal sequence
\[
 S_0=S,\qquad S_{\alpha+1}=\mathsf F(S_\alpha),\qquad
 S_\lambda=\bigcup_{\beta<\lambda}S_\beta
       \quad(\lambda\text{ a nonzero limit ordinal}).
\]
Once equality \(S_\alpha=S_{\alpha+1}\) occurs, write
\(\operatorname{Cl}_{\mathcal F}(S)=S_\alpha\).
This notation is defined on admissible seeds; it does
not assert that a conflicting set has a complete
extension containing it.
\end{hbdefinition}

\begin{hbtheorem}{thm:p18-transfinite-grounded}{proved=annex:p18.transfinite.grounded}{Least complete supersets and the grounded extension}
The displayed sequence is increasing and every stage
is admissible. It reaches equality. If \(A\) is infinite
of cardinality \(\kappa\), equality occurs at some
\(\alpha<\kappa^+\), where \(\kappa^+\) is the next
cardinal viewed as its initial ordinal. If \(A\) has
\(n\) elements, equality occurs by stage \(n\).

The resulting \(\operatorname{Cl}_{\mathcal F}(S)\) is
the least fixed point containing \(S\), and is complete.
In particular
\[
 G_{\mathcal F}=\operatorname{Cl}_{\mathcal F}(\varnothing)
\]
is the unique least fixed point among all fixed points,
including conflicting ones. It is complete, is contained
in every complete extension, and satisfies
\[
 G_{\mathcal F}=\bigcap_{E\in\operatorname{Comp}(\mathcal F)}E.
\]
The closure is extensive, monotone on admissible seeds,
and idempotent on complete extensions.
\end{hbtheorem}

\begin{hbtheorem}{thm:p18-complete-semilattice}{proved=annex:p18.complete.semilattice}{The precise complete-extension order structure}
Every nonempty family \(\mathcal E\) of complete
extensions has a greatest lower bound, given by
\[
 \bigwedge\nolimits_{\operatorname{Comp}}\mathcal E
 =
 \bigcup\{S\in\operatorname{Adm}(\mathcal F):
                           S\subseteq E\text{ for every }E\in\mathcal E\}.
\]
Every nonempty directed family \(\mathcal D\) of
complete extensions has least upper bound
\[
 \bigvee\nolimits_{\operatorname{Comp}}\mathcal D
                  =\operatorname{Cl}_{\mathcal F}\!\left(\bigcup\mathcal D\right).
\]
The empty family's least upper bound is \(G_{\mathcal F}\).
More generally a nonempty family of complete extensions
has a common complete upper bound exactly when its
union is conflict-free; in that case the same closure
formula gives its least upper bound.

Thus complete extensions have the complete-semilattice
structure of Dung's Theorem~25, with directed suprema
as well as the source's increasing-sequence suprema.
They need not form a complete lattice. A meet need
not be the set intersection, and a directed supremum
need not be the set union.
\end{hbtheorem}

\begin{hbdefinition}{def:p18-finitary-framework}{Finitary frameworks and continuity}
The framework is \emph{finitary} if every
\(\operatorname{Att}(a)\) is finite. The set \(A\)
itself may be infinite. Preservation of increasing
countable unions means
\[
 \mathsf F\!\left(\bigcup_{n<\omega}S_n\right)
                  =\bigcup_{n<\omega}\mathsf F(S_n)
                   \quad(S_0\subseteq S_1\subseteq\cdots).
\]
This is the source's \(\omega\)-continuity condition.
\end{hbdefinition}

\begin{hbtheorem}{thm:p18-finitary-continuity}{proved=annex:p18.finitary.continuity}{Finite attacker sets and countable-stage closure}
If the framework is finitary, \(\mathsf F\) preserves
unions of all nonempty directed families of subsets
of \(A\), hence is \(\omega\)-continuous. For every
admissible seed \(S\),
\[
                 \operatorname{Cl}_{\mathcal F}(S)
                      =\bigcup_{n<\omega}\mathsf F^n(S).
\]
In particular this formula computes the grounded set
as a set-theoretic union. It does not imply equality
at a finite stage on an infinite framework.
Finitarity is sufficient and is not necessary for
the displayed continuity.
In a finitary framework the union of a nonempty
directed family of complete extensions is itself
complete, and is their least upper bound.
\end{hbtheorem}

\section{The alternating finite-level hierarchy}

\begin{hbdefinition}{def:p18-alternating-levels}{The no-attack operator and eventual level membership}
Define
\[
 \mathsf N_{\mathcal F}(S)=A\setminus S^+,\qquad
 L_0=A,\qquad L_{n+1}=\mathsf N_{\mathcal F}(L_n).
\]
Write \(\mathsf N\) for this operator in a fixed framework.
An argument is \emph{indefeasible in the finite-level
hierarchy} if some natural number \(m\) satisfies
\(a\in L_n\) for every \(n>m\).
This is Dung's abstract presentation of the
level hierarchy in his comparison with Pollock;
the result below concerns precisely these
sets and this eventual-membership definition
\cite[sec.~4.2, definition~44, pp.~341--342]{Dung1995}.
\end{hbdefinition}

\begin{hbtheorem}{thm:p18-alternating-levels}{proved=annex:p18.alternating.levels}{Alternation, finite-level indefeasibility and grounded semantics}
The operator \(\mathsf N\) reverses inclusion and
\(\mathsf F=\mathsf N\circ\mathsf N\).
For every \(n<\omega\),
\[
 L_{2n}=\mathsf F^n(A),\qquad
 L_{2n+1}=\mathsf F^{n+1}(\varnothing).
\]
The odd levels increase, the even levels decrease,
and every odd level is contained in
\(G_{\mathcal F}\), which is contained in
every even level. The indefeasible set is
\[
 I_{\mathcal F}=\bigcup_{n<\omega}L_{2n+1}
               =\bigcup_{n<\omega}\mathsf F^n(\varnothing)
               \subseteq G_{\mathcal F}.
\]
If the framework is finitary, equality holds.
The inclusion can be strict in a non-finitary
framework. These are the full abstract
claims of Lemmas~45--46 and Theorem~47.
\end{hbtheorem}

\section{Residual frameworks and graph conditions}

\begin{hbdefinition}{def:p18-residual-framework}{Arguments left undecided by a complete extension}
For a complete extension \(E\), put
\[
 B_E=A\setminus(E\cup E^+),\qquad
              \mathcal F\!\upharpoonright B_E=(B_E,R\cap(B_E\times B_E)).
\]
Every use of defense or completeness inside this
induced framework quantifies only over \(B_E\).
\end{hbdefinition}

\begin{hblemma}{thm:p18-residual-lifting}{proved=annex:p18.residual.lifting}{The full residual lifting lemma}
No attack runs in either direction between \(E\)
and \(B_E\). Every member of \(B_E\) has an
attacker in \(B_E\). If \(S\) is admissible
in the residual framework, \(E\cup S\) is
admissible in the original framework.
If \(S\) is complete in the residual framework,
\(E\cup S\) is complete in the original one.
The residual is empty exactly when \(E\) is stable.
\end{hblemma}

\begin{hbdefinition}{def:p18-graph-conditions}{Well-foundedness and the two parity relations}
The framework is \emph{well-founded} when there is
no infinite sequence \(a_0,a_1,\ldots\) with
\((a_{i+1},a_i)\in R\) for every \(i\).
An infinite sequence may repeat vertices.

For \(a,b\in A\), a finite attack walk from \(b\) to \(a\) is a
sequence \(b=v_m,v_{m-1},\ldots,v_0=a\)
with \((v_{i+1},v_i)\in R\). Repeated vertices
are allowed. We say \(b\) \emph{indirectly attacks}
\(a\) if such a walk has odd length, and
\emph{indirectly defends} \(a\) if it has even
length, including zero. Define
\[
 C(b,a)\quad\Longleftrightarrow\quad
 \text{both an odd and an even attack walk run from }b\text{ to }a.
\]
This is Dung's \emph{controversial with respect to}
relation. The framework is \emph{uncontroversial}
when \(C\) is empty. It is \emph{limited
controversial} when there is no infinite sequence
with \(C(a_{i+1},a_i)\) for every \(i\).

It is \emph{coherent} when every preferred
extension is stable, and \emph{relatively
grounded} when
\[
            G_{\mathcal F}
              =\bigcap_{E\in\operatorname{Pref}(\mathcal F)}E.
\]
Preferred extensions have already been proved to
exist. These definitions do not change the
guarded acceptance-query convention of the
finite chapter \cite[definitions~29, 31--32]{Dung1995}.
\end{hbdefinition}

\begin{hbtheorem}{thm:p18-well-founded-extensions}{proved=annex:p18.well.founded.extensions}{Coincidence under well-foundedness}
Every well-founded framework has exactly one
complete extension. It is grounded, preferred
and stable. The result includes infinite
frameworks and the empty framework.
\end{hbtheorem}

\begin{hbdefinition}{def:p18-threat-defense-construction}{The full defense construction}
A \emph{threat} to \(S\) is an argument \(b\)
that attacks some member of \(S\) but is
not attacked by \(S\). A \emph{defense of \(S\)}
is a set that attacks every threat.

Suppose every argument has an attacker, and
fix \(a\in A\) such that \(C(b,a)\) is false
for every \(b\). Define
\[
\begin{aligned}
 E_0&=\{a\},\\
 T_n&=\{b\in A:b\notin E_n^+\land\exists e\in E_n\ (b,e)\in R\},\\
 D_n&=\{d\in A:\exists b\in T_n\ (d,b)\in R\},\\
 E_{n+1}&=E_n\cup D_n,\qquad E_\omega=\bigcup_{n<\omega}E_n.
\end{aligned}
\]
The definition takes all attackers of the threats.
It does not require an inclusion-minimal defense.
\end{hbdefinition}

\begin{hblemma}{thm:p18-noncontroversial-seed}{proved=annex:p18.noncontroversial.seed}{A defended seed and a nonempty complete extension}
Under the stated premises, every member of
\(E_n\) indirectly defends \(a\); every
\(E_n\) is conflict-free; and \(E_\omega\)
is a nonempty admissible set containing
\(a\). Its complete closure therefore
contains \(a\).

Every nonempty limited controversial framework
has a nonempty complete extension. The
nonempty argument-set hypothesis is necessary
for this latter assertion.
\end{hblemma}

\begin{hbtheorem}{thm:p18-coherence}{proved=annex:p18.coherence}{Limited controversy implies coherence and stable existence}
Every limited controversial framework is
coherent and has a stable extension.
Uncontroversial frameworks are limited
controversial, so they also have these
properties \cite[theorem~33(1), corollary~36]{Dung1995}.
\end{hbtheorem}

\begin{hbtheorem}{thm:p18-relative-grounding}{proved=annex:p18.relative.grounding}{Both residual alternatives and relative grounding}
Let the framework be uncontroversial, and
let \(a\in A\setminus(G_{\mathcal F}\cup G_{\mathcal F}^+)\).
There are complete extensions \(E_1,E_2\)
such that \(a\in E_1\) and \(E_2\)
attacks \(a\). Consequently the framework
is relatively grounded
\cite[lemma~35 and theorem~33(2)]{Dung1995}.
\end{hbtheorem}

\section{Infinite boundaries and finite procedures}

\begin{hbdefinition}{def:p18-arbitrary-examples}{Explicit separating frameworks}
The following lists specify all attacks.
\begin{enumerate}
\item In \(\mathcal F_\omega\), the arguments are
      \(g_n,h_n\) for \(n<\omega\), with
      \(g_n\to h_n\) and \(h_n\to g_{n+1}\).
\item In \(\mathcal F_{\omega+1}\), add an argument
      \(z\) to \(\mathcal F_\omega\), and add
      \(h_n\to z\) for every \(n\).
\item In \(\mathcal F_{\rm reverse}\), the arguments
      are the natural numbers, with \(n+1\to n\).
\item In \(\mathcal F_{\rm shareddefender}\), the arguments
      are \(d,z,b_n\) for \(n<\omega\), with
      \(d\to b_n\) and \(b_n\to z\).
\item In \(\mathcal F_{\rm parity}\), use
      \(\{c,b,a\}\) and \(c\to a,c\to b,b\to a\).
\item Dung's initial example is
      \(\mathcal F_{\rm initial}=(\{i_1,i_2,a\},
      \{(i_1,a),(a,i_1),(i_2,a)\})\).
\item In \(\mathcal F_{\rm nomindefense}\), use
      \(a,b_n,x_{k,j}\), for \(n,k,j<\omega\),
      with \(b_n\to a\), \(x_{k,0}\to b_n\)
      whenever \(n\leq k\), and
      \(x_{k,j+1}\to x_{k,j}\).
\end{enumerate}
\end{hbdefinition}

\begin{hbtheorem}{thm:p18-arbitrary-examples}{proved=annex:p18.arbitrary.examples}{What the infinite and graph examples separate}
\begin{enumerate}
\item \(\mathcal F_\omega\) is finitary and well-founded.
      Its finite grounded stages are
      \(G_m=\{g_n:n<m\}\), and its grounded
      extension is \(G_\omega=\{g_n:n<\omega\}\).
      No finite stage is fixed.
\item \(\mathcal F_{\omega+1}\) is well-founded and
      non-finitary. Its finite stages are the same;
      \(G_\omega=\{g_n:n<\omega\}\) is not fixed,
      and \(G_{\omega+1}=G_\omega\cup\{z\}\)
      is grounded and stable. Thus its operator
      fails \(\omega\)-continuity.
\item \(\mathcal F_{\rm reverse}\) is finitary and
      has no directed cycle, but is not well-founded.
      Its complete extensions are the empty set,
      all even numbers, and all odd numbers.
      The latter two are exactly its preferred
      and stable extensions.
\item \(\mathcal F_{\rm shareddefender}\) is non-finitary,
      but its characteristic operator preserves
      every nonempty directed union:
      \(\mathsf F(S)=\{d\}\cup\{z:d\in S\}\).
\item \(\mathcal F_{\rm parity}\) is limited controversial
      and is not uncontroversial: \(C(c,a)\) holds.
      Its unique complete, preferred and stable
      extension is \(\{c\}\).
\item In \(\mathcal F_{\rm initial}\), the grounded
      stages are \(\varnothing,\{i_2\},\{i_1,i_2\}\).
      The final set is the sole complete,
      preferred and stable extension. This
      computes source Examples~3, 8 and 21.
      Its two-cycle shows that well-foundedness
      is not necessary for this coincidence.
      The finite chapter already computes
      the mutual-pair Examples~9 and 22.
\item \(\mathcal F_{\rm nomindefense}\) is uncontroversial
      and every argument has an attacker, but
      \(\{a\}\) has no inclusion-minimal defense.
      The full defense construction above applies.
      Starting at \(a\), its union is the stable
      extension \(\{a\}\cup\{x_{k,2j}:k,j<\omega\}\).
\end{enumerate}
The mutual-pair complete extensions have no
common complete upper bound for the two singletons.
In the finite chapter's shared-argument example,
the meet of its two nonempty complete extensions
is empty although their set intersection is
\(\{d\}\). That shared-argument framework is
limited controversial and coherent but is
not relatively grounded.
In \(\mathcal F_{\omega+1}\), the indefeasible
set \(I_{\mathcal F}\) is \(\{g_n:n<\omega\}\);
the grounded extension also contains \(z\).
\end{hbtheorem}

\begin{hbmethod}{def:p18-finite-graph-procedures}{Exact domains of the finite procedures}
For an explicitly listed finite framework, enumerate
all subsets, compute conflict-freeness and
\(\mathsf F(S)\), and retain the admissible,
complete and stable subsets by their definitions.
Filter the admissible list by inclusion-maximality
to obtain exactly the preferred extensions.
Grounded iteration uses the finite procedure
already proved.
The resulting finite lists decide coherence
by checking whether every preferred extension
is stable, and relative grounding by comparing
the grounded set with the intersection of
the nonempty preferred family.

To compute the controversy relation, replace
each vertex \(a\) by \((a,0),(a,1)\);
an attack \(a\to b\) gives edges
\((a,\epsilon)\to(b,1-\epsilon)\) for
\(\epsilon=0,1\). For each \(b\), calculate
reachability from \((b,0)\), including its
zero-length walk. Put \(C(b,a)\) in the
relation exactly when both \((a,0)\)
and \((a,1)\) are reachable.
\end{hbmethod}

\begin{hbtheorem}{thm:p18-finite-graph-procedures}{proved=annex:p18.finite.graph.procedures}{Termination and the exact finite graph criteria}
The finite procedures terminate and return the
stated extension families and controversy relation.
A finite framework is well-founded exactly
when its attack graph is acyclic. It is
uncontroversial exactly when the computed
\(C\) relation is empty, and limited
controversial exactly when the directed graph
of \(C\) is acyclic.

The acyclicity test uses only the finite
domain hypothesis. On arbitrary infinite
sets, the definitions prohibit infinite
backward sequences, a stronger requirement
than absence of cycles. Transfinite closure
and the countable union for finitary
frameworks are semantic constructions;
an executable procedure on infinite inputs
requires an effective presentation and
its own termination or convergence result.
\end{hbtheorem}

The formal logic-program correspondences, their graph transfers and
Appendices~A--C are supplied in
\cref{thm:p18-program-stable-correspondence,thm:p18-program-well-founded,thm:p18-clause-fitting,thm:p18-program-graph-transfer,thm:p18-program-examples}. The source's Section~4.1 default correspondence
is \cref{thm:p18-default-stable-bijection}; the direct autoepistemic
interface and its exact prerequisite-free scope are
\cref{thm:p18-default-ael-weak,thm:p18-default-ael-prerequisite-free}.
The Section~3 game correspondence is supplied in
\cref{sec:p18-game-argumentation}; the complete Lucas counterexample
is \cref{sec:p18-lucas-no-solution}. The Section~5 meta-interpreter,
including its full signed outputs and infinite-input scope, is
\cref{sec:p18-argumentation-meta-interpreters}.
The broader autoepistemic translations retain their separate
constructions. The present
chapter completes the abstract-framework
semantic and graph-condition chain through
Corollary~36, together with the abstract
alternating-level results through Theorem~47.
